% GENERATED FILE. Edit canonical lessons, then run npm run generate:books.
% canonical-source-hash: fnv1a64:ea8f0ebed200625e
% canonical-lessons: RU-C411-poputchik, RU-C411-poputka, RU-C411-podvozit, RU-C411-vysazhivat, RU-C411-pristyogivatsya

\chapter{Giving someone a lift}
\label{ch:ru-giving-someone-a-lift}
\hlchaptermodality{411}

\begin{chapteropening}
\textbf{By the end of this chapter:} \emph{I can talk about sharing rides: a fellow traveller, a lift in a passing car, giving someone a lift, dropping them off and buckling up.}

Complete the last lesson of chapter 411: I can talk about sharing rides: a fellow traveller, a lift in a passing car, giving someone a lift, dropping them off and buckling up.
\end{chapteropening}

\section[\texorpdfstring{\ru{попутчик} (popútchik)}{попутчик (popútchik)}]{\ru{попутчик} (popútchik) — a fellow traveller}

\label{lesson:RU-C411-poputchik}



\begin{quote}
Before the new one: say the Russian for to knock down, then the Russian for a dent.
\end{quote}

\subsection*{You'll want to know: \ru{попутчик}}
\textbf{\ru{попутчик}} (\emph{popútchik}) — \textquotedblleft{}a fellow traveller\textquotedblright{}.

\begin{quote}
\ru{В} \ru{поезде} \ru{у} \ru{меня} \ru{был} \ru{весёлый} \ru{попутчик}.
\end{quote}

\begin{quote}
\emph{V póyezde u menyá byl vesyólyy popútchik.}
\end{quote}

\begin{quote}
I had a cheerful fellow traveller on the train.
\end{quote}

Feminine \emph{\ru{попутчица}}. Built on \emph{\ru{по} \ru{пути}} — on the way.

\begin{cousinweb}
\emph{\ru{попутчик}} is built on \emph{\ru{попутный}} 'going the same way', from \emph{\ru{по} \ru{пути}} 'on the way'. \emph{\ru{Путь}} is inherited from Proto-Slavic \emph{pǫtĭ} 'a path, a road' (EtymDB).
\end{cousinweb}

\begin{grammarlens}[title={What you've built}]
One more word for telling, explaining and describing at B1.
\end{grammarlens}

\subsection*{Your turn}
\begin{itemize}
\raggedright
  \item \emph{Say it:} \emph{popútchik}
  \item \emph{Say it:} \emph{popútchik}, once more
  \item \emph{Say it:} a sentence of your own with \emph{popútchik}
  \item \emph{Recall:} say the Russian for to knock down, then the Russian for a dent, then say \emph{popútchik} again
\end{itemize}

\begin{tcolorbox}[breakable,colback=teal!4,colframe=teal!35!black,title={Before you move on}]
What does \ru{попутчик} mean? (\textquotedblleft{}A fellow traveller\textquotedblright{}.) Say it once more.
\end{tcolorbox}
\section[\texorpdfstring{\ru{попутка} (popútka)}{попутка (popútka)}]{\ru{попутка} (popútka) — a lift (in a passing car), a hitched ride}

\label{lesson:RU-C411-poputka}



\begin{quote}
Before the new one: say the Russian for a dent, then the Russian for a fellow traveller.
\end{quote}

\subsection*{You'll want to know: \ru{попутка}}
\textbf{\ru{попутка}} (\emph{popútka}) — \textquotedblleft{}a lift (in a passing car), a hitched ride\textquotedblright{}.

\begin{quote}
\ru{До} \ru{деревни} \ru{мы} \ru{ехали} \ru{на} \ru{попутке}.
\end{quote}

\begin{quote}
\emph{Do derévni my yékhali na popútke.}
\end{quote}

\begin{quote}
We got to the village in a passing car.
\end{quote}

Colloquial. \emph{\ru{Ловить} \ru{попутку}} — to hitch a ride.

\begin{cousinweb}
\emph{\ru{попутка}} is a colloquial short form of \emph{\ru{попутная} \ru{машина}} 'a car going the same way', from \emph{\ru{по} \ru{пути}} 'on the way'. \emph{\ru{Путь}} is inherited from Proto-Slavic \emph{pǫtĭ} (EtymDB).
\end{cousinweb}

\begin{grammarlens}[title={What you've built}]
One more word for telling, explaining and describing at B1.
\end{grammarlens}

\subsection*{Your turn}
\begin{itemize}
\raggedright
  \item \emph{Say it:} \emph{popútka}
  \item \emph{Say it:} \emph{popútka}, once more
  \item \emph{Say it:} a sentence of your own with \emph{popútka}
  \item \emph{Recall:} say the Russian for a dent, then the Russian for a fellow traveller, then say \emph{popútka} again
\end{itemize}

\begin{tcolorbox}[breakable,colback=teal!4,colframe=teal!35!black,title={Before you move on}]
What does \ru{попутка} mean? (\textquotedblleft{}A lift (in a passing car), a hitched ride\textquotedblright{}.) Say it once more.
\end{tcolorbox}
\section[\texorpdfstring{\ru{подвозить} (podvozít)}{подвозить (podvozít)}]{\ru{подвозить} (podvozít) — to give (someone) a lift}

\label{lesson:RU-C411-podvozit}



\begin{quote}
Before the new one: say the Russian for a fellow traveller, then the Russian for a lift.
\end{quote}

\subsection*{You'll want to know: \ru{подвозить}}
\textbf{\ru{подвозить}} (\emph{podvozít}) — \textquotedblleft{}to give (someone) a lift\textquotedblright{}.

\begin{quote}
\ru{Каждое} \ru{утро} \ru{сосед} \ru{подвозит} \ru{меня} \ru{до} \ru{работы}.
\end{quote}

\begin{quote}
\emph{Kázhdoye útro soséd podvózit menyá do rabóty.}
\end{quote}

\begin{quote}
Every morning the neighbour gives me a lift to work.
\end{quote}

Imperfective; perfective \emph{\ru{подвезти}} (\emph{\ru{подвезу}, \ru{подвезёшь}}). Present \emph{\ru{подвожу}, \ru{подвозишь}}.

\begin{cousinweb}
\emph{\ru{подвозить}} is built from \emph{\ru{под}-} 'up to' and \emph{\ru{возить}} 'to carry by vehicle'. \emph{\ru{Возить}} is inherited from Proto-Slavic \emph{voziti} 'to cart, to convey' (EtymDB).
\end{cousinweb}

\begin{grammarlens}[title={What you've built}]
One more word for telling, explaining and describing at B1.
\end{grammarlens}

\subsection*{Your turn}
\begin{itemize}
\raggedright
  \item \emph{Say it:} \emph{podvozít}
  \item \emph{Say it:} \emph{podvozít}, once more
  \item \emph{Say it:} a sentence of your own with \emph{podvozít}
  \item \emph{Recall:} say the Russian for a fellow traveller, then the Russian for a lift, then say \emph{podvozít} again
\end{itemize}

\begin{tcolorbox}[breakable,colback=teal!4,colframe=teal!35!black,title={Before you move on}]
What does \ru{подвозить} mean? (\textquotedblleft{}To give (someone) a lift\textquotedblright{}.) Say it once more.
\end{tcolorbox}
\section[\texorpdfstring{\ru{высаживать} (vysázhivat)}{высаживать (vysázhivat)}]{\ru{высаживать} (vysázhivat) — to drop off (a passenger); to plant out}

\label{lesson:RU-C411-vysazhivat}



\begin{quote}
Before the new one: say the Russian for a lift, then the Russian for to give.
\end{quote}

\subsection*{You'll want to know: \ru{высаживать}}
\textbf{\ru{высаживать}} (\emph{vysázhivat}) — \textquotedblleft{}to drop off (a passenger); to plant out\textquotedblright{}.

\begin{quote}
\ru{Такси} \ru{высаживает} \ru{нас} \ru{у} \ru{входа}.
\end{quote}

\begin{quote}
\emph{Taksí vysázhivayet nas u vkhóda.}
\end{quote}

\begin{quote}
The taxi drops us off at the entrance.
\end{quote}

Imperfective; the perfective is \emph{\ru{высадить}} (\emph{\ru{высажу}, \ru{высадишь}}). Also \emph{\ru{высаживать} \ru{цветы}} — to plant out flowers.

\begin{cousinweb}
\emph{\ru{высаживать}} is built from \emph{\ru{вы}-} 'out' and \emph{-\ru{саживать}}, from \emph{\ru{садить}} 'to seat, to plant', the 'make sit' partner of \emph{\ru{сидеть}} 'to sit'. EtymDB lists \emph{\ru{садить}} without a deeper source.
\end{cousinweb}

\begin{grammarlens}[title={What you've built}]
One more word for telling, explaining and describing at B1.
\end{grammarlens}

\subsection*{Your turn}
\begin{itemize}
\raggedright
  \item \emph{Say it:} \emph{vysázhivat}
  \item \emph{Say it:} \emph{vysázhivat}, once more
  \item \emph{Say it:} a sentence of your own with \emph{vysázhivat}
  \item \emph{Recall:} say the Russian for a lift, then the Russian for to give, then say \emph{vysázhivat} again
\end{itemize}

\begin{tcolorbox}[breakable,colback=teal!4,colframe=teal!35!black,title={Before you move on}]
What does \ru{высаживать} mean? (\textquotedblleft{}To drop off (a passenger); to plant out\textquotedblright{}.) Say it once more.
\end{tcolorbox}
\section[\texorpdfstring{\ru{пристёгиваться} (pristyógivatsya)}{пристёгиваться (pristyógivatsya)}]{\ru{пристёгиваться} (pristyógivatsya) — to fasten one's seat belt, to buckle up}

\label{lesson:RU-C411-pristyogivatsya}



\begin{quote}
Before the new one: say the Russian for to give, then the Russian for to drop off.
\end{quote}

\subsection*{You'll want to know: \ru{пристёгиваться}}
\textbf{\ru{пристёгиваться}} (\emph{pristyógivatsya}) — \textquotedblleft{}to fasten one's seat belt, to buckle up\textquotedblright{}.

\begin{quote}
\ru{В} \ru{машине} \ru{мы} \ru{всегда} \ru{пристёгиваемся}.
\end{quote}

\begin{quote}
\emph{V mashíne my vsegdá pristyógivayemsya.}
\end{quote}

\begin{quote}
We always fasten our seat belts in the car.
\end{quote}

Imperfective; perfective \emph{\ru{пристегнуться}}. \emph{\ru{Пристегнуть} \ru{ремень}} — to fasten a belt.

\begin{cousinweb}
\emph{\ru{пристёгиваться}} is built from \emph{\ru{при}-} 'onto', \emph{-\ru{стёгивать}}, which goes with \emph{\ru{стегать}} 'to stitch, to lash', and \emph{-\ru{ся}}. EtymDB does not list the verb, so the deeper origin is left open.
\end{cousinweb}

\begin{grammarlens}[title={What you've built}]
One more word for telling, explaining and describing at B1.
\end{grammarlens}

\subsection*{Your turn}
\begin{itemize}
\raggedright
  \item \emph{Say it:} \emph{pristyógivatsya}
  \item \emph{Say it:} \emph{pristyógivatsya}, once more
  \item \emph{Say it:} a sentence of your own with \emph{pristyógivatsya}
  \item \emph{Recall:} say the Russian for to give, then the Russian for to drop off, then say \emph{pristyógivatsya} again
\end{itemize}

\begin{tcolorbox}[breakable,colback=teal!4,colframe=teal!35!black,title={Before you move on}]
What does \ru{пристёгиваться} mean? (\textquotedblleft{}To fasten one's seat belt, to buckle up\textquotedblright{}.) Say it once more.
\end{tcolorbox}
